\newpage

% =====================================================================
% Załączniki A–P: Architecture Decision Records (skrócone)
% Pełne wersje (format MADR) w repozytorium: docs/adr/
% =====================================================================

\appendix{ADR-0001: Przyjęcie ADR jako narzędzia projektowania i formatu MADR}

\textbf{Status:} zaakceptowany (data: 2026-04-06).

\textbf{Kontekst.} W pracy o tematyce DevOps i SDLC istotne jest udokumentowanie nie tylko stanu końcowego, lecz także uzasadnienia decyzji technicznych. Potrzebny był mechanizm wymuszający podejście „architektura w pierwszej kolejności" oraz przejrzysty ślad decyzyjny na potrzeby oceny akademickiej.

\textbf{Decyzja.} Przyjęto Architecture Decision Records (ADR) w formacie MADR (\textit{Markdown Architectural Decision Records}) jako podstawowe narzędzie sterujące projektem. Każda istotna decyzja jest opisywana w osobnym, wersjonowanym dokumencie.

\textbf{Rozważane alternatywy:} same wiadomości commitów; nieformalna dokumentacja typu wiki.

\textbf{Konsekwencje:}
\begin{itemize}
  \item[$+$] każda zmiana architektoniczna ma osobny, wersjonowany zapis,
  \item[$+$] praca może bezpośrednio cytować ADR-y w uzasadnieniu wyboru rozwiązań,
  \item[$-$] niewielki narzut procesowy --- dokumentacja musi poprzedzać lub towarzyszyć implementacji.
\end{itemize}

\appendix{ADR-0002: Wybór Kubernetes jako głównego środowiska uruchomieniowego}

\textbf{Status:} zaakceptowany (data: 2026-04-12).

\textbf{Kontekst.} Demonstracja praktyk DevOps i cloud-native wymaga standardowej, skalowalnej platformy orkiestracji kontenerów, wspierającej konfigurację deklaratywną, samonaprawę oraz bogaty ekosystem narzędzi CI/CD, GitOps i obserwowalności.

\textbf{Decyzja.} Wybrano Kubernetes (K8s) jako de facto standard aplikacji cloud-native, dający abstrakcje niezbędne dla infrastruktury jako kodu i przepływów GitOps oraz przenośność między środowiskiem lokalnym a chmurą.

\textbf{Rozważane alternatywy:} Docker Compose; rozwiązania serverless/FaaS; tradycyjne maszyny wirtualne.

\textbf{Konsekwencje:}
\begin{itemize}
  \item[$+$] wysoka zgodność z tematem pracy oraz przenośność między chmurami,
  \item[$+$] obsługa automatycznych wdrożeń, skalowania i obserwowalności,
  \item[$-$] większa złożoność i wyższy narzut zasobowy niż w prostszych alternatywach.
\end{itemize}

\appendix{ADR-0003: Wybór ekosystemu Java/JVM dla backendu}

\textbf{Status:} zaakceptowany (data: 2026-04-12).

\textbf{Kontekst.} Backend wymaga dojrzałego, łatwego w utrzymaniu stosu wspierającego nowoczesne wzorce architektoniczne, konteneryzację i obserwowalność, a zarazem wykorzystującego kompetencje autora.

\textbf{Decyzja.} Wybrano ekosystem Java/JVM ze względu na dojrzałość bibliotek i biegłość autora. Aby sprostać ograniczeniom cloud-native (czas startu, zużycie pamięci), projekt celuje w obrazy natywne (GraalVM, kompilacja z wyprzedzeniem).

\textbf{Rozważane alternatywy:} Go; Node.js (TypeScript); Python.

\textbf{Konsekwencje:}
\begin{itemize}
  \item[$+$] maksymalne wykorzystanie kompetencji autora oraz dojrzałych narzędzi DI i testowania,
  \item[$+$] obrazy natywne dają start poniżej sekundy i niski ślad pamięciowy,
  \item[$-$] kompilacja natywna zwiększa złożoność i czas łańcucha CI/CD oraz ogranicza refleksję.
\end{itemize}

\appendix{ADR-0004: Przyjęcie Conventional Commits dla historii zmian}

\textbf{Status:} zaakceptowany (data: 2026-04-21).

\textbf{Kontekst.} Wraz z rozwojem automatyzacji potrzebny był ustandaryzowany, czytelny maszynowo sposób opisu zmian w systemie kontroli wersji, umożliwiający automatyczne generowanie list zmian i wersjonowanie semantyczne.

\textbf{Decyzja.} Przyjęto Conventional Commits --- lekką konwencję nadającą wiadomościom commitów ustrukturyzowaną, czytelną maszynowo formę.

\textbf{Rozważane alternatywy:} commity dowolne; Gitmoji.

\textbf{Konsekwencje:}
\begin{itemize}
  \item[$+$] umożliwia automatyczne określanie przyrostu wersji (semver) i generowanie list zmian,
  \item[$+$] czytelniejsza i przeszukiwalna historia, dobra integracja z CI/CD,
  \item[$-$] wymaga dyscypliny i narzędzi wymuszających format (linting).
\end{itemize}

\appendix{ADR-0005: Strategia konteneryzacji --- Dockerfile jako standard budowania}

\textbf{Status:} zaakceptowany (data: 2026-04-26).

\textbf{Kontekst.} Aplikację Java trzeba w powtarzalny sposób pakować w obraz kontenera; istnieją trzy główne podejścia różniące się przejrzystością, kontrolą i zależnościami narzędziowymi, co wpływa na projekt łańcucha CI/CD oraz walor dydaktyczny pracy.

\textbf{Decyzja.} Wybrano Dockerfile ze względu na pełną przejrzystość procesu, status standardu branżowego oraz możliwość jawnego pokazania dobrych praktyk (budowanie wieloetapowe, użytkownik nie-root, kolejność warstw).

\textbf{Rozważane alternatywy:} Jib (wtyczka Gradle/Maven); Paketo Buildpacks (CNCF).

\textbf{Konsekwencje:}
\begin{itemize}
  \item[$+$] w pełni jawny, czytelny proces budowania --- zgodny z celem dydaktycznym,
  \item[$+$] uniwersalny standard bez zależności od wtyczek; natywna dostępność na runnerach GitHub Actions,
  \item[$-$] ręczne utrzymanie pliku oraz aktualizacja obrazu bazowego pod kątem łatek bezpieczeństwa.
\end{itemize}

\appendix{ADR-0006: OpenTelemetry jako ujednolicony standard obserwowalności}

\textbf{Status:} proponowany (data: 2026-04-26).

\textbf{Kontekst.} Obserwowalność opiera się na trzech filarach (logi, metryki, ślady rozproszone). Bez wspólnego standardu każdy filar wymaga osobnej biblioteki instrumentacji, co wiąże aplikację z konkretnymi dostawcami i utrudnia zmianę zaplecza.

\textbf{Decyzja.} Wybrano OpenTelemetry (CNCF) --- neutralny wobec dostawców standard obejmujący wszystkie trzy filary jednym SDK i Collectorem, z agentem Java instrumentującym aplikację automatycznie, bez zmian w kodzie.

\textbf{Rozważane alternatywy:} osobne biblioteki (Micrometer + SLF4J/Logback + Brave/Zipkin); agent SaaS „wszystko w jednym" (Datadog/Dynatrace).

\textbf{Konsekwencje:}
\begin{itemize}
  \item[$+$] aplikacja odseparowana od konkretnego zaplecza --- zmiana wymaga rekonfiguracji Collectora, a nie kodu,
  \item[$+$] automatyczna instrumentacja (HTTP, JDBC, gRPC) bez zmian w kodzie,
  \item[$-$] wciąż ewoluująca specyfikacja oraz dodatkowy komponent (Collector) do utrzymania.
\end{itemize}

\appendix{ADR-0007: Stos LGTM jako zaplecze obserwowalności}

\textbf{Status:} proponowany (data: 2026-04-26).

\textbf{Kontekst.} Po przyjęciu OpenTelemetry (ADR-0006) należało wybrać konkretne, samodzielnie hostowane i otwarte zaplecza dla logów, metryk i śladów oraz wspólną warstwę wizualizacji, zgodne z Kubernetes i Collectorem OTel.

\textbf{Decyzja.} Wybrano stos LGTM --- Loki (logi), Grafana (wizualizacja), Tempo (ślady) i Prometheus (metryki) --- ponieważ wszystkie komponenty integrują się natywnie przez Grafanę, przyjmują dane z Collectora OTel i wdraża się je oficjalnymi chartami Helm.

\textbf{Rozważane alternatywy:} stos ELK/EFK (Elasticsearch/OpenSearch + Kibana); zarządzany SaaS (Grafana Cloud/Datadog/New Relic).

\textbf{Konsekwencje:}
\begin{itemize}
  \item[$+$] korelacja danych z trzech filarów w jednym interfejsie (Grafana),
  \item[$+$] Loki i Tempo są oszczędne zasobowo --- odpowiednie dla jednowęzłowego klastra akademickiego,
  \item[$-$] cztery osobno zarządzane komponenty; model indeksowania etykiet w Loki wymaga dyscypliny (kardynalność).
\end{itemize}

\appendix{ADR-0008: Wybór Gradle jako systemu budowania monorepozytorium}

\textbf{Status:} zaakceptowany (data: 2026-05-03).

\textbf{Kontekst.} Monorepozytorium zawiera wiele komponentów (backend Java, frontend, infrastruktura). Potrzebny był wydajny system budowania wspierający projekty wielomodułowe, zarządzanie zależnościami oraz buforowanie i zrównoleglenie.

\textbf{Decyzja.} Wybrano Gradle ze względu na najlepszy kompromis między elastycznością a wygodą dla monorepozytorium Java, wydajność (budowy przyrostowe, build cache) oraz dobrą integrację z Quarkusem.

\textbf{Rozważane alternatywy:} Maven (oparty na XML); Bazel (dla bardzo dużych monorepozytoriów).

\textbf{Konsekwencje:}
\begin{itemize}
  \item[$+$] proste zarządzanie zależnościami między modułami i szybkie budowy w GitHub Actions,
  \item[$-$] nieco wyższy próg wejścia niż Maven oraz dodatkowe pliki konfiguracji w katalogu głównym.
\end{itemize}

\appendix{ADR-0009: GitHub jako platforma VCS i CI/CD}

\textbf{Status:} zaakceptowany (data: 2026-05-17).

\textbf{Kontekst.} Projekt wymaga systemu kontroli wersji i platformy CI/CD wspierającej monorepozytorium, integrację z rejestrami obrazów oraz definiowanie łańcuchów w formie kodu (YAML), bezpłatnej dla repozytoriów publicznych i studenckich.

\textbf{Decyzja.} Wybrano GitHub --- standard branżowy z najbogatszym ekosystemem gotowych akcji, natywną integracją z Cosign i Trivy oraz zerowym narzutem infrastrukturalnym; bezpłatny plan wystarcza na potrzeby projektu.

\textbf{Rozważane alternatywy:} GitLab; Bitbucket; Gitea.

\textbf{Konsekwencje:}
\begin{itemize}
  \item[$+$] gotowe akcje (Trivy, Cosign, Docker) oraz integracja Jira z GitHub (pełna identyfikowalność zmian),
  \item[$-$] zależność od zamkniętej platformy chmurowej i limit minut runnerów w planie bezpłatnym.
\end{itemize}

\appendix{ADR-0010: Trivy jako skaner podatności obrazów kontenerów}

\textbf{Status:} zaakceptowany (data: 2026-05-17).

\textbf{Kontekst.} Łańcuch CI musi skanować każdy zbudowany obraz pod kątem znanych podatności (CVE) przed wysłaniem do rejestru (praktyka \textit{Shift Left Security}). Potrzebny był skaner z natywną integracją z GitHub Actions, obejmujący pakiety systemowe i zależności aplikacji oraz przerywający budowę według progu istotności.

\textbf{Decyzja.} Wybrano Trivy (Aqua Security) ze względu na pierwszorzędną akcję GitHub, szersze pokrycie typów pakietów oraz wysoką adopcję w ekosystemie Kubernetes/Cloud Native.

\textbf{Rozważane alternatywy:} Grype (Anchore).

\textbf{Konsekwencje:}
\begin{itemize}
  \item[$+$] zerokonfiguracyjna integracja, status projektu CNCF, wynik SARIF dla zakładki GitHub Security,
  \item[$-$] aktualizacja bazy wymaga dostępu do sieci w CI; sporadyczne fałszywe alarmy przy skanowaniu JVM.
\end{itemize}

\appendix{ADR-0011: Cosign --- bezkluczowe podpisywanie obrazów}

\textbf{Status:} zaakceptowany (data: 2026-05-28).

\textbf{Kontekst.} Po zbudowaniu i przeskanowaniu obrazu łańcuch CI musi poświadczyć jego pochodzenie (ang. \textit{provenance}), bez przechowywania długotrwałych kluczy w repozytorium.

\textbf{Decyzja.} Wybrano Cosign w trybie bezkluczowym (\textit{keyless}, Sigstore): podpis składany jest krótkotrwałym certyfikatem wystawionym na tożsamość OIDC GitHub Actions i zapisywany w publicznym dzienniku przejrzystości Rekor, co umożliwia weryfikację bez dystrybucji klucza publicznego.

\textbf{Rozważane alternatywy:} Cosign ze statyczną parą kluczy; Notary v2 / Docker Content Trust.

\textbf{Konsekwencje:}
\begin{itemize}
  \item[$+$] brak długotrwałych kluczy do rotacji; publicznie weryfikowalne podpisy (dziennik Rekor); zgodność z SLSA,
  \item[$-$] weryfikacja wymaga dostępu sieciowego do Rekor/Fulcio; tryb offline wymaga innego korzenia zaufania.
\end{itemize}

\appendix{ADR-0012: Helm jako menedżer pakietów Kubernetes}

\textbf{Status:} zaakceptowany (data: 2026-06-01).

\textbf{Kontekst.} Manifesty YAML z wartościami zakodowanymi na sztywno (tagi obrazów, przestrzeń nazw, limity) stały się trudne w utrzymaniu przy siedmiu wdrożeniach: tagi wymagały podmiany przez \texttt{sed}, brakowało historii wydań, a wycofanie nieudanego wdrożenia oznaczało ręczne przywracanie manifestów.

\textbf{Decyzja.} Wybrano Helm --- jedyną opcję łączącą parametryzację, wersjonowanie wydań i wycofywanie w jednym narzędziu, będącą standardem w ekosystemie Kubernetes i integrującą się natywnie z ArgoCD/Flux.

\textbf{Rozważane alternatywy:} surowe manifesty z podmianą przez \texttt{sed}; Kustomize.

\textbf{Konsekwencje:}
\begin{itemize}
  \item[$+$] idempotentne wdrożenia, atomowe wycofywanie oraz jeden udokumentowany plik wartości (\texttt{values.yaml}),
  \item[$+$] charty to standardowy mechanizm dystrybucji aplikacji Kubernetes z natywną obsługą w ArgoCD/Flux,
  \item[$-$] rozwlekła składnia szablonów Go; sekrety przekazane przez \texttt{--set} trafiają do historii wydań.
\end{itemize}

\appendix{ADR-0013: ArgoCD jako narzędzie dostawy GitOps}

\textbf{Status:} zaakceptowany (data: 2026-06-09).

\textbf{Kontekst.} Po przyjęciu Helma (ADR-0012) projekt potrzebował kontrolera GitOps domykającego pętlę dostawy: obserwującego repozytorium i uzgadniającego stan klastra. Dominują dwa narzędzia CNCF (ang. \textit{graduated}) --- ArgoCD i Flux.

\textbf{Decyzja.} Wybrano ArgoCD ze względu na pierwszorzędną obsługę chartów Helm z rejestru OCI, obserwowalny interfejs webowy (natychmiast pokazujący stan synchronizacji --- przydatny w demonstracjach i jako materiał ilustracyjny) oraz pozycję najczęściej cytowanego narzędzia GitOps.

\textbf{Rozważane alternatywy:} Flux (GitOps Toolkit, bez wbudowanego interfejsu).

\textbf{Konsekwencje:}
\begin{itemize}
  \item[$+$] wizualizacja drzewa zasobów i różnic Helm w interfejsie; definicje (CRD \texttt{Application}) wersjonowane w Git,
  \item[$-$] większy ślad pamięciowy (około 300 MB) niż Flux; interfejs webowy to dodatkowa powierzchnia ataku.
\end{itemize}

\appendix{ADR-0014: Wybór Vue 3 jako frameworka frontendowego}

\textbf{Status:} zaakceptowany (data: 2026-05-28).

\textbf{Kontekst.} Architektura klient--serwer wymaga warstwy frontendowej --- lekkiej nakładki na REST API, z obsługą TypeScript i dobrą integracją z Vite. Frontend jest elementem pomocniczym, więc liczą się produktywność i prostota.

\textbf{Decyzja.} Wybrano Vue~3 (z Vite i TypeScript) ze względu na najlepszy kompromis między produktywnością, obsługą TypeScript a natywną integracją z Vite dla lekkiej aplikacji jednostronicowej.

\textbf{Rozważane alternatywy:} React (Vite + TypeScript); Angular; Svelte/SvelteKit.

\textbf{Konsekwencje:}
\begin{itemize}
  \item[$+$] Vite jako natywne narzędzie budujące (szybka iteracja, małe paczki produkcyjne); komponenty jednoplikowe czytelne w prezentacji akademickiej,
  \item[$-$] mniejszy ekosystem bibliotek UI niż React oraz mniejsza popularność w zastosowaniach korporacyjnych niż Angular/React.
\end{itemize}

\appendix{ADR-0015: k3d jako lokalne środowisko Kubernetes}

\textbf{Status:} zaakceptowany (data: 2026-06-08).

\textbf{Kontekst.} Projekt wymaga lokalnego klastra Kubernetes do rozwoju, testów integracyjnych i demonstracji (ADR-0002), szybkiego w utworzeniu i udostępniającego Ingress pod adresem \texttt{http://localhost} bez obejść. Początkowo użyto minikube, lecz w trybie sterownika Docker wymagał on uruchamiania \texttt{minikube tunnel} w osobnym terminalu, a czas startu był długi.

\textbf{Decyzja.} Wybrano k3d --- cienką nakładkę na k3s uruchamiającą klaster w kontenerach Docker. Mapowanie portów na etapie tworzenia klastra usuwa konieczność tunelowania, start trwa poniżej 30 sekund, a wbudowany Traefik pełni rolę kontrolera Ingress.

\textbf{Rozważane alternatywy:} minikube; kind; Kubernetes w Docker Desktop.

\textbf{Konsekwencje:}
\begin{itemize}
  \item[$+$] Ingress dostępny pod \texttt{http://localhost} bez tunelu; szybki start (około 20 s); proste klastry wielowęzłowe (\texttt{--agents N}),
  \item[$-$] k3s pomija część rzadko używanych API; domyślny Ingress to Traefik (manifesty muszą używać \texttt{ingressClassName: traefik}).
\end{itemize}

\appendix{ADR-0016: Sealed Secrets do zarządzania sekretami Kubernetes}

\textbf{Status:} zaakceptowany (data: 2026-06-11).

\textbf{Kontekst.} Model GitOps (ADR-0013) wymaga przechowywania wszystkich zasobów w repozytorium Git, lecz natywny \texttt{Secret} koduje dane jedynie w Base64 --- jego commit jest równoznaczny z ujawnieniem haseł. Potrzebny był sposób utrzymania sekretów w postaci zaszyfrowanej w repozytorium, uzgadnialny przez ArgoCD bez zewnętrznego systemu i bez ręcznych kroków poza pętlą GitOps.

\textbf{Decyzja.} Wybrano Bitnami Sealed Secrets --- kontroler działający w klastrze szyfruje \texttt{Secret} kluczem publicznym do zasobu \texttt{SealedSecret} bezpiecznego do commitowania, a klucz prywatny nigdy nie opuszcza klastra. W konsekwencji chart Helm przestał renderować jakikolwiek zasób \texttt{Secret}: dane niewrażliwe trafiają do \texttt{ConfigMap}, a hasła pochodzą wyłącznie z odszyfrowanych \texttt{SealedSecret}.

\textbf{Rozważane alternatywy:} External Secrets Operator (synchronizacja z zewnętrznego magazynu); HashiCorp Vault; SOPS (szyfrowanie plików kluczem age/GPG).

\textbf{Konsekwencje:}
\begin{itemize}
  \item[$+$] sekrety zaszyfrowane w Git i uzgadnialne przez ArgoCD bez zewnętrznych zależności; brak haseł w jawnym tekście w \texttt{values.yaml},
  \item[$-$] klucz prywatny kontrolera stanowi pojedynczy punkt zaufania (wymaga kopii zapasowej), a rotacja sekretu oznacza ręczne ponowne uszczelnienie i commit.
\end{itemize}

% =====================================================================
% Załącznik Q: skorowidz kart narzędzi
% Pełne karty znajdują się w treści rozdziału 3; źródła w docs/tech-card/
% =====================================================================

\newpage
\appendix{Karty narzędzi --- skorowidz}

Pełne karty narzędzi (opis, fragment konfiguracji, dobre nawyki, rola w projekcie) znajdują się w treści części praktycznej (rozdział~3) jako wstawki przy odpowiednich etapach pętli DevOps. Wersje źródłowe w formacie Markdown są przechowywane w repozytorium projektu w katalogu \texttt{docs/tech-card/}. Tabela~\ref{tab:tech-cards} zestawia wszystkie karty wraz z etapem ich omówienia oraz ścieżką pliku źródłowego.

\begin{table}[H]
  \centering
  \renewcommand{\thetable}{\arabic{table}}%
  \caption{Skorowidz kart narzędzi}
  \label{tab:tech-cards}
  \begin{tabular}{lll}
    \hline
    \textbf{Narzędzie} & \textbf{Etap DevOps} & \textbf{Plik źródłowy} \\
    \hline
    Quarkus                  & Code    & \texttt{docs/tech-card/quarkus.md} \\
    Vue 3                    & Code    & \texttt{docs/tech-card/vue.md} \\
    Docker                   & Build   & \texttt{docs/tech-card/docker.md} \\
    Helm                     & Release & \texttt{docs/tech-card/helm.md} \\
    Kubernetes               & Deploy  & \texttt{docs/tech-card/kubernetes.md} \\
    k3d                      & Deploy  & \texttt{docs/tech-card/k3d.md} \\
    ArgoCD                   & Deploy  & \texttt{docs/tech-card/argocd.md} \\
    Sealed Secrets           & Deploy  & \texttt{docs/tech-card/sealed-secrets.md} \\
    OpenTelemetry Collector  & Monitor & \texttt{docs/tech-card/opentelemetry-collector.md} \\
    Prometheus               & Monitor & \texttt{docs/tech-card/prometheus.md} \\
    Loki                     & Monitor & \texttt{docs/tech-card/loki.md} \\
    Tempo                    & Monitor & \texttt{docs/tech-card/tempo.md} \\
    Grafana                  & Monitor & \texttt{docs/tech-card/grafana.md} \\
    \hline
  \end{tabular}
\end{table}
